\chapter*{General Introduction}
\addcontentsline{toc}{chapter}{General Introduction}
\adjustmtc
\thispagestyle{MyStyle}

\lettrine{R}{egulatory} compliance in international trade has traditionally
relied on manual document review: an analyst reads a supplier's mill
certificate or carbon-intensity report, cross-references it against a
country's published default emission values, and manually computes a
liability. As the European Union's \ac{CBAM} enters its definitive,
annual-declaration phase in 2026, this manual burden falls on every EU
importer of steel and aluminium and, symmetrically, creates new reporting
pressure on the non-EU exporters whose shipment documents now feed
directly into a foreign compliance calculation. \acp{LLM} offer an
obvious lever to automate the reading and reconciliation of such
documents, but a
compliance number is exactly the kind of output where an LLM's
well-documented tendency to hallucinate plausible-looking
values~\cite{ji2023hallucination} cannot be tolerated.

This project resolves that tension architecturally, rather than by simply
hoping the model behaves. A \textbf{deterministic, LLM-free calculation
engine} computes every compliance number. Around it sits a
\textbf{genuinely agentic tool-use loop}: a locally hosted, open-weight
LLM that decides for itself which extraction, lookup, and validation tool
to call, in what order, until it is confident enough to write a result or
escalates to a human instead. The two halves run on separate machines,
communicating over a private network, and are integrated into a working
Odoo~17 module --- so what comes out the other end is a declaration
workflow a real compliance team could open and use, not a research
prototype left in isolation.

This report is organised into three chapters. Chapter~\ref{chap:context}
establishes the regulatory context, states the problem, and reviews the
literature on agentic AI, tool use, and hallucination risk that motivates
the system's design. Chapter~\ref{chap:design} presents the system
analysis and design: requirements, actors and use cases, the sequence and
class diagrams, the design philosophy, and the system architecture.
Chapter~\ref{chap:implementation} covers the implementation --- the
calculation engine, the agentic pipeline, and the empirical evaluation
results --- and closes with a candid discussion of the system's
limitations. A general conclusion summarises the contribution and
outlines future work.
